\section{Sensor Calibration}\label{sec:calibration}

Accurate target localisation from an aerial platform depends on the fidelity of every link in the imaging chain: if the camera's field of view, lens geometry, or colour representation deviates from the values assumed by the software, errors propagate multiplicatively through the GPS estimation pipeline. This section documents the four calibration campaigns conducted during development, each of which corrected a specific source of systematic error.

% ─────────────────────────────────────────────────────────────────────
\subsection{Field-of-View Calibration}\label{sec:cal:fov}

The ground footprint of a pinhole camera at altitude $h$ is governed by
\begin{equation}\label{eq:footprint-calib}
    W_{\text{ground}} = \frac{w_{\text{sensor}} \cdot h}{f}
\end{equation}
where $w_{\text{sensor}}$ is the physical sensor width and $f$ is the effective focal length. This relationship is used twice in the system: first by \texttt{planning.py} to compute the scan-strip spacing (with a 20\% overlap margin), and second by the GPS estimation module to convert pixel offsets from the image centre into metric displacements on the ground. An error of $\Delta f / f$ in the focal length therefore produces an identical fractional error in every GPS estimate and in the strip spacing, meaning the search pattern either leaves uncovered gaps or wastes flight time on excessive overlap.

The IMX296 global-shutter sensor has a nominal pixel pitch of \SI{3.45}{\micro\metre} and an active area width of \SI{5.02}{\milli\metre}. The lens datasheet quoted a focal length of \SI{7.0}{\milli\metre}, which was used as the initial default in \texttt{config.py}. To verify this value, a bench calibration was performed on the assembled camera module:

\begin{enumerate}
    \item The camera was mounted vertically at a measured height of \SI{1.000 \pm 0.005}{\metre} above a flat surface, using the \texttt{capture\_training.py} MJPEG stream for live preview.
    \item A tape measure was placed horizontally across the full width of the frame.
    \item The visible extent was read directly from the tape as \SI{0.92 \pm 0.01}{\metre}.
\end{enumerate}

Substituting into Equation~\ref{eq:footprint-calib} and solving for $f$:
\begin{equation}\label{eq:fov-result}
    f = \frac{w_{\text{sensor}} \cdot h}{W_{\text{ground}}} = \frac{5.02 \times 1.00}{0.92} = \SI{5.46}{\milli\metre}
\end{equation}

The calibrated value is 22\% lower than the datasheet nominal. The corresponding horizontal field of view is
\begin{equation}
    \theta_{\text{HFOV}} = 2 \arctan\!\left(\frac{w_{\text{sensor}}}{2f}\right) = 2 \arctan\!\left(\frac{5.02}{2 \times 5.46}\right) \approx 49.3^{\circ}
\end{equation}

Had the uncalibrated value of \SI{7.0}{\milli\metre} been retained, every ground-distance calculation would have been scaled by a factor of $5.46/7.0 = 0.78$---i.e.\ the system would have \emph{underestimated} the camera footprint by 22\%, producing GPS target estimates displaced from the true position by the same fraction of the pixel offset. At a search altitude of \SI{35}{\metre}, the ground footprint is \SI{32.2}{\metre} (calibrated) versus \SI{25.1}{\metre} (uncalibrated); a detection at the frame edge would yield a GPS error of $\approx$\SI{1.5}{\metre} purely from the focal length mismatch. The strip spacing in the search pattern would also have been 22\% too narrow, wasting approximately four additional scan lines over the survey area.

% ─────────────────────────────────────────────────────────────────────
\subsection{Lens Distortion Calibration}\label{sec:cal:lens}

Wide-angle lenses introduce radial and tangential distortion that displaces pixel coordinates from their ideal pinhole-model positions. For a target near the frame edge, uncorrected barrel distortion can shift the detected bounding-box centre by several pixels, degrading the pixel-to-GPS conversion. The standard Brown--Conrady model~\cite{brown1966decentering} parameterises radial distortion as
\begin{equation}\label{eq:distortion}
    r' = r\bigl(1 + k_1 r^2 + k_2 r^4 + k_3 r^6\bigr)
\end{equation}
where $r$ is the normalised radial distance from the principal point and $k_1, k_2, k_3$ are the radial distortion coefficients. Tangential terms ($p_1, p_2$) account for misalignment between the lens and sensor planes.

Calibration was performed using a printed checkerboard target (calib.io pattern, $14 \times 9$ squares at \SI{28}{\milli\metre} pitch, yielding $13 \times 8$ inner corners). Approximately 20 images were captured at varying orientations using \texttt{tests/calibration/lens\_calibrate.py}, which calls \texttt{cv2.findChessboardCornersSB} (the sector-based detector, more robust to lighting variation) with a fallback to the standard \texttt{cv2.findChessboardCorners} with adaptive thresholding and corner normalisation flags. Sub-pixel refinement was applied with \texttt{cv2.cornerSubPix} using a $(5,5)$ search window and termination criteria of 30 iterations or $\epsilon = 0.001$.

The resulting reprojection RMS was \textbf{0.399\,px}, indicating an excellent fit. The five distortion coefficients and $3\times3$ camera matrix were saved to \texttt{calibration\_data.npz}. At runtime, \texttt{vision.py} loads this file during initialisation and precomputes a pair of remap lookup tables via \texttt{cv2.initUndistortRectifyMap}:

\begin{verbatim}
new_mtx, _ = cv2.getOptimalNewCameraMatrix(
    mtx, dist, (cam_w, cam_h), 0, (cam_w, cam_h))
map1, map2 = cv2.initUndistortRectifyMap(
    mtx, dist, None, new_mtx, (cam_w, cam_h), cv2.CV_16SC2)
\end{verbatim}

Every frame is then undistorted with a single call to \texttt{cv2.remap(frame, map1, map2, cv2.INTER\_LINEAR)} before being passed to the detection model. The precomputed integer remap (\texttt{CV\_16SC2}) adds only \SI{1.5}{\milli\second} per frame---negligible compared to the \SI{207}{\milli\second} TFLite inference time---so the correction is effectively free. All downstream modules (passive observer, ground station, autonomous mission) benefit transparently because undistortion is applied inside \texttt{detect\_in\_image()}.

% ─────────────────────────────────────────────────────────────────────
\subsection{Camera Colour-Space Discovery}\label{sec:cal:colour}

The IMX296 sensor is configured via \texttt{picamera2} with the \texttt{RGB888} pixel format, which the library documents as returning 8-bit RGB data. However, during integration testing on the Raspberry Pi~5, colour reproduction appeared incorrect: the detection model (trained on BGR images, the OpenCV default) produced anomalously low confidence scores, and visual inspection of the stream showed a red/blue channel swap.

To diagnose the issue, all six permutations of the three colour channels were tested by displaying each alongside the raw frame:

\begin{verbatim}
for perm in [(0,1,2), (0,2,1), (1,0,2), (1,2,0), (2,0,1), (2,1,0)]:
    cv2.imshow(str(perm), frame[:,:,list(perm)])
\end{verbatim}

The identity permutation $(0,1,2)$---i.e.\ treating the data as BGR without any conversion---produced correct colours. This confirms that the IMX296 driver delivers data in \textbf{BGR} byte order despite the \texttt{RGB888} format label. Applying the na\"ive \texttt{cv2.cvtColor(frame, cv2.COLOR\_RGB2BGR)} conversion that the label would suggest actually \emph{swaps} the channels into an incorrect order, resulting in the red and blue planes being exchanged.

The fix was to remove the incorrect \texttt{cvtColor} call entirely. The practical impact is twofold. First, the YOLO model---which expects BGR input---now receives correctly ordered data, restoring detection confidence to its trained level ($>0.95$ on the bench-test dummy). Second, the MJPEG stream served to the ground-station browser displays natural colours without post-processing, reducing per-frame latency. This discovery was documented as a permanent warning in the project documentation and \texttt{config.py} to prevent future regressions.

% ─────────────────────────────────────────────────────────────────────
\subsection{DJI Video FOV Calibration}\label{sec:cal:dji}

Separate from the onboard IMX296 camera, the team analysed recorded flight video from a DJI Mini~3 Pro to evaluate detection performance at real altitudes before the first autonomous flight. The DJI footage was centre-cropped from $3840\times2160$ to $1456\times1088$ to match the onboard camera's aspect ratio, but the crop changes the effective field of view, which must be known to convert pixel detections into metric ground coordinates.

A calibration tool (\texttt{tools/fov\_calibrate\_video.py}) was developed. The operator pauses the video at frames where the mannequin is visible and clicks the top and bottom of the dummy; the tool computes the apparent height in pixels and, using the known real height of \SI{1.8}{\metre} and the SRT-encoded altitude, solves for the focal length in pixels:

\begin{equation}\label{eq:dji-focal}
    f_{\text{px}} = \frac{h_{\text{px}} \cdot h_{\text{alt}}}{h_{\text{real}}}
\end{equation}

Nine samples were collected across altitudes from \SI{15}{\metre} to \SI{50}{\metre}, yielding a mean focal length of $1416 \pm 41$\,px. The corresponding horizontal FOV for the $1456$-pixel crop width is

\begin{equation}
    \theta_{\text{HFOV}} = 2\arctan\!\left(\frac{1456}{2 \times 1416}\right) = 54.4^{\circ}
\end{equation}

\noindent This DJI video HFOV ($54.4^{\circ}$) is distinct from the onboard Pi IMX296 camera HFOV ($49.3^{\circ}$, Section~\ref{sec:fov-results}); the two values correspond to different cameras with different optics.

As a consistency check, the dummy height was back-calculated at each sample altitude using the calibrated focal length; all nine measurements fell within \SIrange{1.7}{1.9}{\metre}, confirming the calibration against the known \SI{1.8}{\metre} ground truth. This calibrated value replaced a hardcoded $73^{\circ}$ assumption, which would have underestimated the focal length by 34\% and corrupted every GPS estimate derived from the DJI analysis footage.

% ─────────────────────────────────────────────────────────────────────
\subsection{Calibration Impact Summary}\label{sec:cal:summary}

Table~\ref{tab:calibration} summarises each calibration, its method, and the quantified impact on system accuracy.

\begin{table}[ht]
\centering
\caption{Sensor calibration summary. The ``Error if uncalibrated'' column quantifies the systematic bias that would have affected every measurement had the default or assumed value been retained.}
\label{tab:calibration}
\renewcommand{\arraystretch}{1.25}
\begin{tabular}{@{}p{2.8cm}cccp{4.2cm}@{}}
\toprule
\textbf{Calibration} & \textbf{Before} & \textbf{After} & \textbf{Method} & \textbf{Error if uncalibrated} \\
\midrule
Focal length (IMX296) & \SI{7.0}{\milli\metre} & \SI{5.46}{\milli\metre} & Tape at \SI{1}{\metre} & 22\% GPS estimate bias; 4 wasted scan lines \\
Lens distortion & None & RMS 0.399\,px & Checkerboard $13{\times}8$ & Edge-pixel shift of several px; \SI{<0.5}{\metre} GPS error at \SI{35}{\metre} \\
Colour space & RGB assumed & BGR actual & 6-permutation test & Model confidence $\downarrow$; colour-dependent features corrupted \\
DJI video FOV & $73^{\circ}$ assumed & $54.4^{\circ}$ measured & Click calibration, 9 samples & 34\% error in video-derived GPS estimates \\
\bottomrule
\end{tabular}
\end{table}

The four calibrations are complementary: the FOV calibration corrects the dominant scale factor in the pinhole model, lens undistortion corrects the nonlinear residual, the colour-space fix ensures the detection model receives data in its trained domain, and the DJI FOV calibration enables accurate offline analysis of flight footage. Together, they reduce the systematic component of the target localisation error to the level set by GPS receiver noise ($\sim$\SI{2}{\metre} CEP) and atmospheric effects, rather than being dominated by avoidable modelling errors.

% ─────────────────────────────────────────────────────────────────────
\subsection{Pre-Flight Calibration Checklist}\label{sec:cal:checklist}

Before every flight, the following calibration and readiness checks must be completed in order. Table~\ref{tab:preflight-cal} presents the checklist with acceptance criteria and the consequence of skipping each item.

\begin{table}[ht]
\centering
\caption{Pre-flight calibration checklist. Each item must meet the stated acceptance criterion before proceeding to flight. Items marked $\ast$ are optional but recommended.}
\label{tab:preflight-cal}
\renewcommand{\arraystretch}{1.3}
\begin{tabularx}{\linewidth}{@{}c >{\raggedright}p{3.6cm} >{\raggedright}p{3.8cm} X@{}}
\toprule
\textbf{\#} & \textbf{Check} & \textbf{Acceptance Criterion} & \textbf{If Failed} \\
\midrule
1 & FOV calibration & Visible width at \SI{1}{\metre} = \SI{92 \pm 3}{\centi\metre} & Recompute \texttt{FOCAL\_LENGTH\_MM}; GPS estimates biased by $\Delta f / f$ \\
2 & Camera colour order & Stream shows natural colours (BGR, no \texttt{cvtColor}) & Detection confidence drops; retrained model receives wrong channel order \\
3 & GPS fix quality & Fix type $\geq 3$, satellites $\geq 10$, HDOP $< 1.5$ & Do not arm; position estimates unreliable \\
4 & Detection confidence & Point camera at mannequin, confidence $> 0.4$ & Check model file, lighting, focus \\
5\textsuperscript{$\ast$} & Lens undistortion & \texttt{calibration\_data.npz} valid ($>$\SI{1}{\kilo\byte}) & Delete corrupt file; edge-pixel error $\leq$\SI{0.3}{\metre} at \SI{35}{\metre} \\
6 & Altitude zero-check & Baro reads \SI{0}{\metre} on ground ($\pm$\SI{0.5}{\metre}) & Recalibrate barometer in Mission Planner \\
7 & Inference speed & $<$\SI{250}{\milli\second} (TFLite) or $<$\SI{100}{\milli\second} (NCNN) & Check CPU throttling, model file size \\
\bottomrule
\end{tabularx}
\end{table}

% ─────────────────────────────────────────────────────────────────────
\subsection{GPS Error Measurement Procedure}\label{sec:cal:gps-error}

The GPS receiver introduces an irreducible random error that sets the floor of target localisation accuracy. Quantifying this error before flight allows the operator to judge whether conditions are adequate and provides a baseline against which CV-derived estimates are compared.

\paragraph{Equipment.} Drone with GPS antenna, known reference position (e.g.\ the designated Take-Off Location at 51.423406\textdegree\,N, $-$2.671446\textdegree\,W).

\paragraph{Procedure.}
\begin{enumerate}
    \item Place the drone at the known reference GPS position with a clear sky view. Allow the GPS receiver to acquire a stable fix (fix type $\geq 3$, $\geq 10$ satellites).
    \item Run the automated drift measurement script:
    \begin{verbatim}
    python tests/day_1_experiments/gps_drift.py
    \end{verbatim}
    The script records GPS latitude, longitude, altitude, satellite count, and HDOP at \SI{1}{\hertz} for \SI{60}{\second}, writing results to a timestamped CSV file.
    \item Compute the following metrics from the recorded data:
    \begin{itemize}
        \item \textbf{CEP50}: the radius of a circle centred on the mean position that contains 50\% of readings. This is the standard measure of GPS accuracy.
        \item \textbf{CEP95}: the radius containing 95\% of readings.
        \item \textbf{Max deviation}: the largest single-sample displacement from the mean.
    \end{itemize}
    \item Record the satellite count and HDOP at the time of measurement for traceability.
\end{enumerate}

\paragraph{Acceptance criteria.} CEP50 $< \SI{3}{\metre}$ with $\geq 10$ satellites and HDOP $< 1.5$. With a Here~3+ GNSS receiver in open sky, typical values are CEP50 $\approx \SIrange{1.5}{2.5}{\metre}$.

\paragraph{Impact.} GPS position error adds directly to target estimation error---it cannot be reduced by improving the vision pipeline. If CEP50 $= \SI{2}{\metre}$, the target estimate can never be better than \SI{2}{\metre} even with perfect FOV calibration and zero lens distortion. For a comprehensive decomposition of all error sources, see the error budget in Table~\ref{tab:error-budget} (Appendix~\ref{sec:error-budget}).

% ─────────────────────────────────────────────────────────────────────
\subsection{Altitude Error Measurement Procedure}\label{sec:cal:alt-error}

ArduCopter reports altitude relative to the arming point using a barometric pressure sensor. Because barometric pressure varies with temperature, humidity, and weather, the reported altitude drifts over time. Since altitude enters directly into the ground sample distance (GSD) calculation---$\text{GSD} = w_s \cdot h / (f \cdot W_{\text{px}})$---any altitude error produces a proportional error in every pixel-to-metre conversion.

\paragraph{Procedure.}
\begin{enumerate}
    \item Place the drone on flat, level ground.
    \item Arm the flight controller (propellers removed for bench testing) and note the initial altitude reading. It should read \SI{0.0}{\metre} ($\pm$\SI{0.1}{\metre}).
    \item Wait \SI{5}{\minute}, monitoring the altitude readout via:
    \begin{verbatim}
    python tests/hardware/gps_test.py
    \end{verbatim}
    \item Record the altitude at the end of the waiting period.
    \item If the drift exceeds \SI{0.5}{\metre}, recalibrate the barometer in Mission Planner (\textit{Initial Setup} $\rightarrow$ \textit{Mandatory} $\rightarrow$ \textit{Accel Calibration}).
\end{enumerate}

\paragraph{Quantified impact.} Table~\ref{tab:alt-error-impact} shows the position error introduced by altitude measurement error at the frame edge, where the pixel-to-GPS conversion is most sensitive.

\begin{table}[ht]
\centering
\caption{Position error at the frame edge caused by barometric altitude error, computed as $\Delta x = (\Delta h / h) \times (W_{\text{ground}} / 2)$ where $W_{\text{ground}}$ is the ground footprint width at the nominal altitude.}
\label{tab:alt-error-impact}
\renewcommand{\arraystretch}{1.15}
\begin{tabular}{@{}S[table-format=2.0] S[table-format=1.1] S[table-format=1.1] S[table-format=1.1]@{}}
\toprule
{\textbf{Altitude (m)}} & {\textbf{$\Delta h = \SI{0.5}{\metre}$}} & {\textbf{$\Delta h = \SI{1.0}{\metre}$}} & {\textbf{$\Delta h = \SI{2.0}{\metre}$}} \\
\midrule
10 & 0.2 & 0.5 & 0.9 \\
20 & 0.2 & 0.5 & 0.9 \\
35 & 0.2 & 0.5 & 0.9 \\
50 & 0.2 & 0.5 & 0.9 \\
\bottomrule
\end{tabular}

\smallskip
\footnotesize Note: the fractional error $\Delta h / h$ varies with altitude, but the absolute edge error converges because the frame footprint also scales with $h$. The dominant effect is on GSD accuracy and, consequently, on the strip spacing in the search pattern.
\end{table}

\paragraph{Acceptance criterion.} Altitude drift $< \SI{0.5}{\metre}$ over \SI{5}{\minute} on the ground. If conditions are marginal (hot day, rapidly changing pressure), re-arm immediately before takeoff to reset the altitude reference.

% ─────────────────────────────────────────────────────────────────────
\subsection{Error Budget Cross-Reference}\label{sec:cal:error-xref}

The calibrations documented in this appendix address only the \emph{systematic} (correctable) components of the target localisation error. The full error budget---including random GPS noise, GPS timing lag, attitude-induced projection shift, and motion blur---is presented in Table~\ref{tab:error-budget} (Appendix~\ref{sec:error-budget}). Table~\ref{tab:cal-vs-budget} maps each calibration in this appendix to the error source it mitigates, providing a unified view of which calibrations reduce which budget entries.

\begin{table}[ht]
\centering
\caption{Mapping between calibration procedures (this appendix) and the error budget entries in Table~\ref{tab:error-budget}. ``Residual after cal.'' is the remaining error contribution at \SI{35}{\metre} altitude after the calibration is applied.}
\label{tab:cal-vs-budget}
\renewcommand{\arraystretch}{1.2}
\begin{tabular}{@{}p{3.2cm} p{3.5cm} S[table-format=1.1] p{3.8cm}@{}}
\toprule
\textbf{Calibration} & \textbf{Error Budget Entry} & {\textbf{Residual (m)}} & \textbf{Section} \\
\midrule
FOV / focal length & GSD scale factor & 0.0 & \S\ref{sec:cal:fov} \\
Lens undistortion & Pixel coordinate shift & 0.02 & \S\ref{sec:cal:lens} \\
Colour-space fix & Detection confidence & {---} & \S\ref{sec:cal:colour} \\
GPS error measurement & Receiver position noise & 2.0 & \S\ref{sec:cal:gps-error} \\
Altitude error check & Barometric drift & 0.5 & \S\ref{sec:cal:alt-error} \\
\midrule
\multicolumn{2}{@{}l}{\textbf{Combined RSS (calibrated)}} & 2.1 & \\
\multicolumn{2}{@{}l}{\textbf{Combined RSS (uncalibrated)}} & 4.8 & \\
\bottomrule
\end{tabular}
\end{table}

\noindent Calibration reduces the systematic error contribution from \SI{4.8}{\metre} to \SI{2.1}{\metre} RSS, bringing the total localisation uncertainty to a level dominated by the irreducible GPS receiver noise rather than avoidable modelling errors. The \SI{2.1}{\metre} calibrated RSS is consistent with the CEP50 $\approx \SI{2.3}{\metre}$ measured from DJI flight video analysis (Section~\ref{sec:cal:dji}).
